\section{连续正则化流（CNF）的基本框架}

\subsection{三个核心概念}

Flow Matching 的数学框架围绕三个核心概念展开，它们之间存在紧密的相互联系：

\begin{importantbox}{Flow Matching 的三大核心概念}
\begin{enumerate}
    \item \textbf{流映射}（Flow Map）$\phi_t(x)$：从起始点 $x$ 出发，沿轨迹到达时刻 $t$ 的位置
    \item \textbf{向量场}（Vector Field）$v_t(x)$：定义 ODE 中每个时空点处的速度方向
    \item \textbf{概率密度路径}（Probability Density Path）$p_t(x)$：时刻 $t$ 处数据点的概率分布
\end{enumerate}
这三者之间的关系贯穿整个 Flow Matching 理论。
\end{importantbox}

\subsection{流映射（Flow Map）}

流映射 $\phi_t(x_0)$ 表示从初始点 $x_0$（来自参考分布 $p_0$）出发，经过时间 $t$ 后到达的位置。形式化定义为：

$$\phi_t(x_0) = x_0 + \int_0^t v_s(\phi_s(x_0)) \, ds$$

其中各符号的含义为：
\begin{itemize}
    \item $\phi_t(x_0)$：从 $x_0$ 出发，在时刻 $t$ 的轨迹位置
    \item $v_s(\cdot)$：时刻 $s$ 处的向量场（速度场）
    \item $\phi_0(x_0) = x_0$：边界条件，$t=0$ 时轨迹起点就是 $x_0$ 本身
    \item $\phi_1(x_0) \sim p_1$：边界条件，$t=1$ 时应到达数据分布的样本
\end{itemize}

\begin{knowledgebox}{ODE 与流映射的关系}
流映射本质上是以下常微分方程（ODE）的解：
$$\frac{d\phi_t(x_0)}{dt} = v_t(\phi_t(x_0)), \quad \phi_0(x_0) = x_0$$
这里\textbf{没有随机性}——一旦确定了起始点 $x_0$ 和向量场 $v_t$，整条轨迹就完全确定了。这也是"Ordinary" Differential Equation 中"Ordinary"的含义——相对于 SDE 中的随机项。
\end{knowledgebox}

有时为了简化符号，我们会用 $x_t$ 来表示 $\phi_t(x_0)$，即从 $x_0$ 出发在时刻 $t$ 的轨迹点。

\subsection{向量场（Vector Field）}

向量场 $v_t(x)$ 定义了每个时空点 $(t, x)$ 处的"速度"。它是流映射的时间导数：

$$v_t(\phi_t(x_0)) = \frac{d\phi_t(x_0)}{dt}$$

向量场编码了从噪声到数据的映射信息。如果我们能找到一个好的向量场，使得从 $p_0$ 采样的点能被映射到 $p_1$（数据分布），那么我们就得到了一个生成模型。

\subsection{概率密度路径（Probability Density Path）}

概率密度路径 $p_t(x)$ 描述了时刻 $t$ 处数据点 $x$ 的概率密度。它满足以下性质：
\begin{itemize}
    \item $p_0 = \mathcal{N}(0, I)$（参考分布，通常为标准高斯）
    \item $p_1 \approx p_{\text{data}}$（数据分布）
    \item 中间时刻 $p_t$ 是从 $p_0$ 到 $p_1$ 的插值
\end{itemize}

\begin{knowledgebox}{连续性方程（Continuity Equation）}
向量场和概率密度路径之间的关系由\textbf{连续性方程}给出：
$$\frac{\partial p_t(x)}{\partial t} + \nabla \cdot (p_t(x) v_t(x)) = 0$$
这是 Fokker-Planck 方程在无扩散项（即 ODE 而非 SDE）情况下的简化形式。它建立了向量场和概率路径之间的精确数学联系。
\end{knowledgebox}

\subsection{三者之间的关系}

在扩散模型中，我们的推导路径是：
$$\text{概率密度路径 } p_t \;\longrightarrow\; \text{SDE/ODE 形式} \;\longrightarrow\; \text{向量场 } v_t \;\longrightarrow\; \text{流映射 } \phi_t$$

而在 Flow Matching 中，我们走的是\textbf{相反的路径}：
$$\text{向量场 } v_t \;\longrightarrow\; \text{流映射 } \phi_t \;\longrightarrow\; \text{概率密度路径 } p_t$$

也就是说，我们先定义向量场，再由它推导出流映射和概率路径。

\subsection{对数概率的计算}

向量场还可以用来计算数据点的对数概率。对于 $x_1 \sim p_1$，其对数似然可以通过以下方式计算：

$$\log p_1(x_1) = \log p_0(x_0) - \int_0^1 \nabla \cdot v_t(\phi_t(x_0)) \, dt$$

其中各项的含义为：
\begin{itemize}
    \item $\log p_0(x_0)$：起始点在参考分布下的对数概率（高斯分布，已知解析形式）
    \item $\nabla \cdot v_t$：向量场的散度（divergence）
    \item 积分项：沿轨迹的体积变化累积
\end{itemize}

这意味着，如果我们能找到一个使 $\log p_1(x_1)$ 最大化的向量场，就能训练出一个好的生成模型。

\subsection{本章小结}

连续正则化流（CNF）的框架由三个核心概念构成：流映射、向量场和概率密度路径。它们之间通过 ODE 和连续性方程相互联系。与扩散模型不同，CNF 直接从向量场出发定义生成过程，提供了更直接的数学框架。

